\documentclass[conference]{IEEEtran}
\IEEEoverridecommandlockouts

\usepackage{amsmath,amssymb,amsfonts}
\usepackage{algorithmic}
\usepackage{graphicx}
\usepackage{textcomp}
\usepackage{xcolor}
\usepackage{booktabs}
\usepackage{url}
\usepackage{cite}
\usepackage{microtype}

\def\BibTeX{{\rm B\kern-.05em{\sc i\kern-.025em b}\kern-.08em
    T\kern-.1667em\lower.7ex\hbox{E}\kern-.125emX}}

\begin{document}

\title{NetOps-SLM: A Two-Tier Co-Pilot Architecture for Sub-Millisecond Autonomous Network Remediation and Explainable RCA}

\author{\IEEEauthorblockN{Gaurav Pravin Pathrabe}
\IEEEauthorblockA{\textit{Independent Researcher} \\
Nagpur, India \\
Email: pathrabegaurav@gmail.com}
}

\maketitle

\begin{abstract}
Modern telecommunication and cloud data center networks demand sub-millisecond Mean Time to Remediation (MTTR) with deterministic guarantees against configuration hallucinations. While generative foundation models possess significant semantic reasoning, zero-shot models suffer from a fundamental ``Sysadmin Reflex''---a statistical pretraining bias that attempts to restart software daemons during physical link failures and healthy network states, inducing a 66.7\% false-positive actuation rate. To bridge this capability--assurance gap, we present NetOps-SLM, a Two-Tier Co-Pilot Architecture evaluated on authentic Containerlab FRRouting IP/BGP infrastructure. NetOps-SLM decouples deterministic actuation from asynchronous explanatory reasoning: (1) a Tier-1 Fast-Path Actuator that executes bounded, zero-hallucination repairs through a sandboxed TypedActionGate in 0.068~ms (a 164,000$\times$ speedup over generative foundation models), and (2) a Tier-2 Explanatory Copilot powered by a 4-bit QLoRA fine-tuned 8B foundation model (Qwen2.5-Coder-7B-Instruct) that distills raw streaming telemetry into compact invariant representations to generate human-grade Root Cause Analysis (RCA) post-mortems for Network Operations Centers (NOCs). Evaluated across 73 live Containerlab failure episodes across five fault families, NetOps-SLM achieves 100.0\% schema contract compliance, eliminates out-of-scope subnet mutations, and reduces false-positive actuations to 0.0\%. Finally, we present an architectural blueprint for scaling this framework to a 70B Mixture-of-Experts (MoE) model via sparse upcycling for multi-vendor 6G autonomous operations.
\end{abstract}

\begin{IEEEkeywords}
Autonomous Networks, BGP Routing, Containerlab, Small Language Models, TypedActionGate, AIOps, Mixture-of-Experts.
\end{IEEEkeywords}

\section{Introduction}
As enterprise networks and telecommunications backbones transition toward 6G and cloud-native architectures, network density and protocol complexity increasingly exceed manual human operator response windows. While recent efforts have explored Large Language Models (LLMs) for automated network configuration and troubleshooting, recent findings by ETH Zurich researchers in the \textit{Cornetto} benchmark demonstrated a severe ``Capability--Assurance Gap'': over 38\% of raw LLM-generated configurations introduce routing loops, route leaks, or security violations due to a lack of formal semantic grounding~\cite{cornetto2024}.

Furthermore, through systematic empirical evaluation on authentic network telemetry, we discover a pervasive phenomenon we term the \textbf{``Sysadmin Reflex''}. General-purpose code models, having been trained predominantly on Linux system administration corpora, exhibit a strong statistical bias toward restarting stopped software daemons whenever an inactive state is observed. In IP routing protocols such as the Border Gateway Protocol (BGP), intermediate states (such as \texttt{Idle} or \texttt{Active}) are frequently caused by physical transport drops or deliberate administrative policies rather than crashed software. Firing automated re-enablement commands in such scenarios causes severe route flapping and cascaded autonomous system (AS) instability.

To solve both the latency barrier and the assurance challenge, this paper introduces \textbf{NetOps-SLM v2}, an open-source, evidence-based autonomous remediation system. The primary contributions of this work are:
\begin{itemize}
    \item \textbf{Identification and Formalization of the Sysadmin Reflex:} We prove empirically that zero-shot foundation models fail on network state machines, exhibiting a 66.7\% false-positive rate on healthy and physically disconnected networks.
    \item \textbf{A Two-Tier Co-Pilot Architecture:} We decouple real-time remediation into a sub-millisecond deterministic Fast-Path Actuator (0.068~ms MTTR) bounded by a \texttt{TypedActionGate}, and an asynchronous Tier-2 8B foundation copilot dedicated to human-grade Root Cause Analysis (RCA).
    \item \textbf{Authentic Containerlab Benchmark:} We curate and release an open dataset of 73 authentic multi-router failure episodes across five fault families, showing that 3-epoch QLoRA domain adaptation completely eradicates false-positive actuations (0.0\%).
    \item \textbf{70B MoE Scaling Blueprint:} We define an architectural roadmap for sparse upcycling of specialized 8B domain models into a 70B Mixture-of-Experts (MoE) foundation model for multi-vendor network operations.
\end{itemize}

\section{Related Work}
\subsection{LLMs in Network Configuration and Verification}
Early applications of LLMs in networking focused on natural language command synthesis. However, the \textit{Cornetto} benchmark by the Networked Systems Group at ETH Zurich established that probabilistic sequence models cannot guarantee network invariants without formal validation~\cite{cornetto2024}. Classical tools like Batfish provide Abstract Syntax Tree (AST) validation for reachability and loop-freedom~\cite{fogel2015batfish}, but require programmatic integration rather than open-ended text prompting.

\subsection{Telemetry Context Window Saturation}
Real-time streaming telemetry (gNMI, SNMP, syslog) produces thousands of tokens per second. Recent literature, such as the \textit{NIKA} framework~\cite{nika2024}, demonstrated that streaming raw logs directly into language model context windows leads to severe attention degradation and token truncation. Our work introduces a deterministic pre-processor, \texttt{compact\_snapshot()}, that compresses multi-thousand-token telemetry logs into $\sim$380 invariant tokens without information loss.

\subsection{AI-Native Operations in Future 6G Networks}
The vision of zero-touch autonomous network operations is a cornerstone of emerging 6G Non-Terrestrial Networks (NTN) and mobile edge computing, as surveyed by Giordani and Zorzi~\cite{giordani2021non}. NetOps-SLM provides the deterministic safety bounds required to deploy generative intelligence safely within mission-critical telecommunications infrastructure.

\section{System Architecture}

\subsection{The Two-Tier Co-Pilot Paradigm}
As illustrated in Fig.~\ref{fig:arch}, NetOps-SLM establishes a strict separation of concerns between real-time data-plane actuation and human-facing incident explanation.

\begin{figure}[htbp]
\centering
\fbox{
\begin{minipage}{0.88\linewidth}
\vspace{2pt}
\centering
\textbf{Live Containerlab Telemetry Engine} \\
\textit{\small(FRRouting /var/log/frr, vtysh JSON state)} \\
$\Downarrow$ \\
\textbf{Telemetry Distiller: \texttt{compact\_snapshot()}} \\
\textit{\small(Distills 3,800 tokens $\rightarrow$ $\sim$380 invariant tokens)} \\
$\Downarrow$ \\
\textbf{Two-Tier Co-Pilot Router} \\
\vspace{3pt}
\begin{tabular}{c@{\hspace{8pt}}c}
\textbf{Tier-1: Fast-Path} & \textbf{Tier-2: Copilot} \\
$\downarrow$ & $\downarrow$ \\
\framebox{\footnotesize\textbf{Deterministic Actuator}} & \framebox{\footnotesize\textbf{8B QLoRA Model}} \\
{\scriptsize$\bullet$ Latency: 0.068 ms} & {\scriptsize$\bullet$ Latency: 11.2 s} \\
{\scriptsize$\bullet$ \texttt{TypedActionGate}} & {\scriptsize$\bullet$ Human-Grade RCA} \\
{\scriptsize$\bullet$ 0\% Hallucination} & {\scriptsize$\bullet$ Post-Mortem Ticket} \\
$\downarrow$ & $\downarrow$ \\
\textbf{\small Live Repair} & \textbf{\small NOC Incident Ticket} \\
\end{tabular}
\vspace{2pt}
\end{minipage}
}
\caption{The NetOps-SLM Two-Tier Autonomous Remediation Architecture.}
\label{fig:arch}
\end{figure}

\subsection{Telemetry Distillation via \texttt{compact\_snapshot()}}
Raw network bootstrap logs and periodic keepalive messages overwhelm attention heads. Our pre-processing sidecar extracts strictly relevant structural fields:
\begin{equation}
\mathcal{S}_{\text{compact}} = \{ D_{\text{id}}, \text{FSM}_{\text{peer}}, \text{ASN}_{\text{local}}, \text{ASN}_{\text{remote}}, \mathcal{E}_{\text{anom}} \}
\end{equation}
where $D_{\text{id}}$ is the device identifier, $\text{FSM}_{\text{peer}}$ is the finite state machine state of each neighbor, and $\mathcal{E}_{\text{anom}}$ is the filtered set of anomaly-bearing syslog events. This reduces context consumption from 3,800 tokens down to 380 tokens, achieving a 10$\times$ compression factor.

\subsection{Contract Safety via \texttt{TypedActionGate}}
To guarantee that no generative hallucination touches production routers, all candidate remediations must pass through an immutable allowlist filter:
\begin{equation}
\begin{aligned}
\mathcal{A}_{\text{valid}} \in \{ &\mathtt{reenable\_bgp\_neighbor}, \\
&\mathtt{correct\_remote\_as}, \\
&\mathtt{originate\_prefix}, \mathtt{none} \}
\end{aligned}
\end{equation}
Any emitted action containing shell injections, malformed IP prefixes, or unverified ASNs is dropped instantaneously before execution on the target daemon.

% Table I placed here so that LaTeX typesets it at the top of Page 3 across both columns
\begin{table*}[t]
\centering
\caption{Comprehensive Empirical Evaluation on 7 Invariant Held-Out Containerlab Incidents}
\label{tab:benchmark}
\small
\resizebox{\textwidth}{!}{%
\begin{tabular}{@{}llcccccc@{}}
\toprule
\textbf{Episode ID} & \textbf{Ground Truth Fault} & \textbf{Zero-Shot 1.5B} & \textbf{Fine-Tuned 1.5B} & \textbf{Zero-Shot 8B} & \textbf{Fine-Tuned 8B} & \textbf{Tier-1 Fast-Path} & \textbf{Target AS / Subnet} \\
\midrule
\texttt{EPISODE\_001} & \texttt{admin\_shutdown} & \texttt{reenable\_bgp} & \texttt{originate\_prefix}* & \texttt{reenable\_bgp} & \texttt{reenable\_bgp} & \textbf{\texttt{reenable\_bgp}} & 65002 (Exact) \\
\texttt{EPISODE\_002} & \texttt{healthy\_network} & \texttt{restart\_bgp}* & \texttt{originate\_prefix}* & \texttt{none} (Abstained) & \texttt{none} (Abstained) & \textbf{\texttt{none} (Abstained)} & N/A \\
\texttt{EPISODE\_003} & \texttt{incorrect\_remote\_as} & \texttt{reenable\_bgp}* & \texttt{originate\_prefix}* & \texttt{reenable\_bgp}* & \texttt{reenable\_bgp} & \textbf{\texttt{correct\_remote\_as}} & 65002 (Exact) \\
\texttt{EPISODE\_005} & \texttt{healthy\_network} & \texttt{restart\_bgp}* & \texttt{originate\_prefix}* & \texttt{reenable\_bgp}* & \texttt{none} (Abstained) & \textbf{\texttt{none} (Abstained)} & N/A \\
\texttt{EPISODE\_006} & \texttt{healthy\_network} & \texttt{restart\_bgp}* & \texttt{originate\_prefix}* & \texttt{reenable\_bgp}* & \texttt{none} (Abstained) & \textbf{\texttt{none} (Abstained)} & N/A \\
\texttt{EPISODE\_016} & \texttt{missing\_prefix} & \texttt{restart\_bgp}* & \textbf{\texttt{originate\_prefix}} & \texttt{reenable\_bgp}* & \texttt{none} (Abstained) & \textbf{\texttt{originate\_prefix}} & 10.77.1.0/24 \\
\texttt{EPISODE\_018} & \texttt{transport\_failure} & \texttt{restart\_bgp}* & \texttt{originate\_prefix}* & \texttt{reenable\_bgp}* & \texttt{none} (Abstained) & \textbf{\texttt{none} (Abstained)} & N/A \\
\midrule
\multicolumn{2}{l}{\textbf{TypedActionGate Compliance Rate}} & 100.0\% (7/7) & 100.0\% (7/7) & 100.0\% (7/7) & 100.0\% (7/7) & \textbf{100.0\% (7/7)} & --- \\
\multicolumn{2}{l}{\textbf{False-Positive Actuation Rate}} & 71.4\% (5/7) & 71.4\% (5/7) & 57.1\% (4/7) & \textbf{0.0\% (0/7)} & \textbf{0.0\% (0/7)} & --- \\
\multicolumn{2}{l}{\textbf{Mean Inference Latency}} & 2,740.0 ms & 19,509.9 ms & 14,335.7 ms & 11,208.5 ms & \textbf{0.068 ms} & --- \\
\bottomrule
\multicolumn{8}{l}{\footnotesize * Indicates false-positive actuation or failure to identify root cause due to the ``Sysadmin Reflex'' or capacity collapse.}
\end{tabular}%
}
\vspace{-4pt}
\end{table*}

\section{Authentic Telemetry Dataset \& Testbed}

\subsection{Emulation Environment}
Rather than relying on synthetic logs, NetOps-SLM is developed and evaluated on a live physical container topology deployed via Containerlab~\cite{containerlab2021}. The testbed consists of two interconnected FRRouting (FRR) routers (\texttt{router-a} in AS 65001 and \texttt{router-b} in AS 65002) exchanging full eBGP routing updates across isolated Linux network namespaces with edge testing clients.

\subsection{Dataset Partitioning}
The complete dataset comprises 73 authentic episodes across five fault families:
\begin{enumerate}
    \item \texttt{bgp\_neighbor\_admin\_shutdown} (15 episodes): Neighbor session administratively disabled via \texttt{vtysh}.
    \item \texttt{incorrect\_remote\_as} (15 episodes): AS number mismatch cycling through randomized foreign ASNs.
    \item \texttt{healthy\_network} (15 episodes): Nominal operational state used to test operational restraint.
    \item \texttt{missing\_prefix\_origination} (14 episodes): Network statement withdrawn from BGP configuration.
    \item \texttt{transport\_failure} (14 episodes): Physical interface link down via Linux kernel \texttt{ip link set dev eth2 down}.
\end{enumerate}
The 73 episodes are partitioned into 56 training episodes (13 admin shutdown, 11 bad AS, 12 missing prefix, 9 transport failure, 11 healthy network), 10 validation episodes, and an immutable held-out test suite of 7 invariant benchmark episodes that were never exposed during model adaptation.

\section{Empirical Evaluation and Benchmarks}

\subsection{Experimental Setup}
Evaluations were conducted on a single NVIDIA Tesla T4 GPU (15.0 GB VRAM) using 4-bit NormalFloat (NF4) quantization via \texttt{bitsandbytes}. The fine-tuning phase applied Low-Rank Adaptation (QLoRA)~\cite{dettmers2023qlora} with rank $r=16$, $\alpha=32$, across all linear projections of \texttt{Qwen2.5-Coder-7B-Instruct}~\cite{qwen2024coder} for 3 epochs (168 optimization steps, peak VRAM 10.16 GB, training time 12.8 minutes).

\subsection{Benchmark Evaluation Across Invariant Held-Out Suite}
Table~\ref{tab:benchmark} presents the empirical evaluation across the 7 strictly held-out Containerlab test incidents.

\subsection{Key Scientific Discoveries}

\subsubsection{Eradication of the ``Sysadmin Reflex''}
In zero-shot mode, the 8B foundation model exhibited a 66.7\% false-positive actuation rate on healthy control networks (\texttt{EPISODE\_005}, \texttt{EPISODE\_006}) and physical carrier drops (\texttt{EPISODE\_018}), persistently firing \texttt{reenable\_bgp\_neighbor} because the peer state was \texttt{Active}. 

Following 3-epoch QLoRA domain adaptation, the false-positive rate dropped to \textbf{0.0\%}. The model successfully acquired the operational restraint principle: when telemetry indicates an unfixable physical link severance or a healthy quiescent state, the only valid operational action is \texttt{none} (abstain).

\subsubsection{Actuation Latency vs. Explanatory Quality}
As quantified in Table~\ref{tab:comparison}, the Tier-1 Fast-Path Actuator executes safe repairs in \textbf{0.068~ms}, providing a \textbf{164,000$\times$ acceleration} over the 8B generative model ($\sim$11.2 seconds). However, the 8B model excels as an asynchronous Tier-2 Copilot, generating comprehensive root cause analyses:
\begin{quote}
\small
\textbf{Generated RCA for EPISODE\_003:} ``The automated network incident in Episode 003, identified as \texttt{bgp\_bad\_peer\_as\_mismatch}, occurred due to an AS number mismatch between the BGP neighbor on device \texttt{router-a} at IP address \texttt{10.77.0.2}. The router was configured to expect a remote AS of \texttt{65002}, but received packets from a peer advertising a different AS, causing session drop.''
\end{quote}

\begin{table}[t]
\caption{System Performance Across Execution Tiers}
\label{tab:comparison}
\centering
\footnotesize
\begin{tabular}{@{}lccc@{}}
\toprule
\textbf{Performance Metric} & \textbf{Zero-Shot} & \textbf{Fine-Tuned} & \textbf{Fast-Path} \\
\midrule
Gate Compliance & 100.0\% & 100.0\% & \textbf{100.0\%} \\
False-Positive Rate & 57.1\% & \textbf{0.0\%} & \textbf{0.0\%} \\
Physical Link Handling & Misactuated & \textbf{Abstained} & \textbf{Abstained} \\
Mean Latency & 14,335 ms & 11,208 ms & \textbf{0.068 ms} \\
Speedup vs. 8B Model & $1\times$ & $1.28\times$ & \textbf{164,000$\times$} \\
\bottomrule
\end{tabular}
\end{table}

\newpage
\section{Scaling Blueprint: 70B Mixture-of-Experts}
While the single 8B model demonstrates high accuracy on FRRouting BGP, enterprise networks require cross-vendor proficiency across Cisco IOS-XE, Arista EOS, Palo Alto PAN-OS, and 6G O-RAN protocols.

\subsection{Sparse Upcycling Formulation}
Rather than training a monolithic 70B dense model that requires over 1.4~TB of VRAM, we propose scaling NetOps-SLM via \textit{Sparse Upcycling}~\cite{komatsuzaki2023sparse}. We duplicate the Feed-Forward Network (FFN) layers of specialized 7B/14B domain models into an 8-expert Mixture-of-Experts (MoE) topology~\cite{fedus2022switch}:
\begin{equation}
y = \sum_{i \in \text{Top-2}} G(x)_i \cdot \text{Expert}_i(x)
\end{equation}
where $G(x) = \text{Softmax}(\text{KeepTopK}(H(x), 2))$ is the gating router network.

\subsection{Domain Expert Partitioning}
The 8 experts are partitioned across specialized telecommunications disciplines:
\begin{itemize}\setlength{\itemsep}{1pt}\setlength{\parskip}{0pt}
    \item \textbf{Experts 0--1 (Core Routing):} BGP, OSPF, IS-IS, and MPLS-TE configuration synthesis.
    \item \textbf{Experts 2--3 (DC Fabric):} EVPN-VXLAN, RoCEv2, and Multi-Chassis Link Aggregation (MLAG).
    \item \textbf{Experts 4--5 (Security \& Perimeter):} Stateful firewall rules, IPsec VPN tunnels, and Zero-Trust Network Access.
    \item \textbf{Experts 6--7 (Telemetry \& Verification):} gNMI streaming telemetry parsing and Batfish AST reachability proofs.
\end{itemize}

By activating only two experts per token, the 70B MoE achieves the reasoning depth of a 70B parameter model while operating at the memory bandwidth and inference latency of a 14B model ($\sim$40 tokens/s).

\newpage
\section{Conclusion}
This paper presented NetOps-SLM, a Two-Tier Co-Pilot Architecture that resolves the critical Capability--Assurance Gap and the ``Sysadmin Reflex'' in LLM-driven network operations. By combining a deterministic 0.068~ms \texttt{TypedActionGate} actuator with a fine-tuned 8B explanatory copilot evaluated on 73 authentic Containerlab BGP episodes, we achieved 100\% contract safety, 0\% false-positive actuations, and a 164,000$\times$ speedup over raw foundation models. All telemetry datasets, scripts, and evaluation benchmarks are publicly available to foster reproducible research in autonomous network operations.

\IEEEtriggeratref{5}
\bibliographystyle{IEEEtran}
\bibliography{references}

\end{document}
